\documentclass{ximera}




% MATH -----------------------------------------------------------
\newcommand{\nats}{\mathbb N}
\newcommand{\ints}{\mathbb Z}
\newcommand{\rats}{\mathbb Q}
\newcommand{\reals}{\mathbb R}
\newcommand{\complex}{\mathbb C}
\newcommand{\powerset}{\mathscr P}
\newcommand{\Laplace}[1]{\ensuremath{\mathcal{L}{\left[#1\right]}}}
\newcommand{\InvLap}[1]{\ensuremath{\mathcal{L}^{-1}{\left[#1\right]}}}

%-------------------------------------------------------------

\pgfplotsset{compat=1.18}
\begin{document}
\begin{abstract}
 \end{abstract}

    \title{Class Discussion Activity 12}
    
    \maketitle

\section*{Exercises}



\begin{enumerate}[label=(\arabic*)]


\item $\displaystyle \mathcal{L}(t^3e^{-2t})=\ldots$

\begin{enumerate}[label=(\alph*)]
\item $\dfrac{6}{s^4}+2$
\item $\dfrac{6}{(s+2)^4}$  % ***
\item $\dfrac{6}{(s-2)^4}$
\item $-\dfrac{2}{s^4+1}$
\item $\dfrac{6}{s^3(s+2)}$ 
\item None of these
\end{enumerate}

% (b) with options through (f)

\item Find the value of $\mathcal{L}((t-1)^2e^{-t})$ at $s=1$.  


% 2/(s+1)^3 -2 * 1/(s+1)^2 +1/(s+1) at s=1 is 1/4-1/2+1/2 = 0.25



\item Let $F(s)=\mathcal{L}(f)$ where $f(t)=\left\{
     \begin{array}{lr}
       4-t^2 & \mbox{ if }0\leq t<2\\
       \\
       0 &  \mbox{ if }t\geq 2
     \end{array}\right.$  Determine $F(3)$ to the nearest thousandth.

     % 1.261

     % 4-t^2 +U(t-2)[t^2-4] ----> 4/s - 2/s^3 +e^(-2s)(L(t^2+4t)) ----> 4/s - 2/s^3 +e^(-2s)(2/s^3+4/s^2)

    % F(3)=4/3-2/27+e^(-6)*(2/27+4/9)\approx 1.261 ... CHECKS OUT WITH DIRECT INTEGRAL CALCULATION




\item  One can show that $\mathcal{L}\left(\sin{(\omega t)}\right)$ is $\displaystyle \frac{\omega}{s^2+\omega^2}$.   Using this, and the first shifting property, find the value of the Laplace transform $\mathcal{L}\left(e^{-3t}\sin{\left(\frac{\pi t}{4}\right)}\right)$ at $s=2$.  Round to the nearest hundredth.

% (0.25pi)/[(s+3)^2+0.0625pi^2] at s=2 is  (0.25pi)/[25+0.0625pi^2]\approx 0.03

\newpage

\item Suppose $f(t)$ is piecewise continuous on $[0,\infty)$ and is of exponential order $s_0\geq 0$.  For $s>s_0$, let $\displaystyle F(s)=\mathcal{L}(f(t))=\int_0^{\infty} f(t)e^{-st}\,\mathrm{d}t$.   Which of the following is $\displaystyle \frac{dF}{ds}$ for $s>s_0$? 


    Hint:  If there exists $G(s,t)$ with continuous 2nd order partial derivatives for $s>s_0$ and $t\geq 0$ satisfying $G_t=f(t)e^{-st}$ for $t\geq 0$ then for $T>0$, $\displaystyle \int_0^{T} f(t)e^{-st}\,\mathrm{d}t=G(s,T)-G(s,0)$ and thus \\$\displaystyle \dfrac{d}{ds}\displaystyle \int_0^{T} f(t)e^{-st}\,\mathrm{d}t=G_s(s,T)-G_s(s,0)=\int_0^T G_{st}\,\mathrm{d}t=\int_0^T G_{ts}\,\mathrm{d}t$.


\begin{enumerate}[label=(\alph*)]
\item $\displaystyle s\mathcal{L}(tf(t))$
\item $\displaystyle -s\mathcal{L}(tf(t))$
\item $\displaystyle \mathcal{L}(tf(t))$
\item $\displaystyle -\mathcal{L}(tf(t))$ % ***
\item $\displaystyle -\mathcal{L}(f(t))$
\item None of these
\end{enumerate}

% (d) with options through (f)


\item  Use the result of the previous exercise, and that $\displaystyle \mathcal{L}\left(\sin{(\omega t)}\right)=\frac{\omega}{s^2+\omega^2}$, to determine the value of the Laplace transform of $t\sin{(2t)}$ at $s=1$.  


    % - d/(ds) (2/(s^2+4))= 4s/(s^2+4)^2.  At s=1, that's 4/25=0.16

\item Suppose $f(t)$ is continuous and of exponential order $s_0>0$.  Using integration by parts, determine which expression below is $\displaystyle \Laplace{\int_0^t f(x)\;dx}$ for $s>s_0$.

% \frac{1}{s}\Laplace{f}

\begin{enumerate}[label=(\alph*)]
\item $\displaystyle \dfrac{1}{s}\Laplace{f}$  % ***
\item $\displaystyle s\Laplace{f}-f(0)$
\item $\displaystyle s\Laplace{f}+f(0)$
\item $\displaystyle -\dfrac{1}{s}\Laplace{f}$
\item $\displaystyle s-\dfrac{1}{s}\Laplace{f}$
\item None of these
\end{enumerate}

% (a) with options through (f)

\newpage

\item Determine the value of the inverse Laplace transform of $\displaystyle F(s)=\frac{s-1}{s^2+2s+5}$ at $t=1$.  Round to the nearest hundredth.

% -0.49


\item Evaluate $\displaystyle \mathcal{L}^{-1}\left(\dfrac{28s^2+2}{16s^4+s^2}\dfrac{}{}\right)$ at $t=2\pi$.  Round to the nearest hundredth.

% 11.57


\item Evaluate $\displaystyle \mathcal{L}^{-1}\left(\frac{s^4-3s^3+18s^2+81}{s^6+18s^4+81s^2}\right)$ at $t=1$ to the nearest hundredth.

    % Ans: f(t) = t-\frac{1}{2}t\sin{(3t)} ---> Laplace: 1/s^2 +0.5 d/(ds) (3/(s^2+9))= 1/s^2 - 3s/(s^2+9)^2

  %  f(1)=1-0.5sin(3)\approx 0.93


\end{enumerate}



\end{document}
